\documentclass[a4paper,12pt]{article}

\usepackage[norwegian]{babel}
\usepackage[utf8]{inputenc}
\usepackage{array}
\usepackage{pdflscape}
\usepackage{geometry}
\usepackage{fancyhdr}
\geometry{margin=2cm}

\begin{document}

% ------------------ FORSIDE (VERTIKAL) ------------------



% ------------------ LANDSKAP STARTER HER ------------------

\begin{landscape}

% ------------------ ØKT 1 ------------------
\subsection*{Undervisningsøkt 1 – Ekskursjon 4 timer}

\begin{tabular}{|p{5cm}|p{6cm}|p{8cm}|p{6cm}|}
\hline
\textbf{Kompetansemål} & \textbf{Læringsmål} & \textbf{Arbeidsmåter} & \textbf{Vurdering} \\
\hline
Naturfag: Å kunne regne i naturfag er å kunne innhente, bearbeide og framstille relevant tallmateriale (...)\newline

Kroppsøving: I kroppsøving handlar det tverrfaglege temaet berekraftig utvikling om naturopplevingar med vekt på trygg og berekraftig ferdsel (...)\newline

bruke kart og digitale verktøy på ein måte som sikrar trygg ferdsel for seg sjølv og for andre\newline


& 
Elevene skal kunne: måle pH-verdier i ulike vannkilder\newline

trives og være trygge på tur i nærmiljøet\newline

finne fram i naturen på egenhånd ved hjelp av kart eller skilting.
& 
Begynner med en kort introduksjon eller samtale om sur nedbør og dens effekter på naturen.\newline

Elevene samler inn pH-verdier fra elv, bekk, vann, myr, osv.\newline

Diskuterer observasjoner og samarbeider i grupper.
& 
Korrekt gjennomføring av målinger\newline

Samarbeid og deltakelse\newline

Forberedelse (riktige klær, mat, osv.)
\\
\hline
\end{tabular}

\newpage

% ------------------ ØKT 2 ------------------
\subsection*{Undervisningsøkt 2 – Analyse 2 timer}

\begin{tabular}{|p{5cm}|p{6cm}|p{8cm}|p{6cm}|}
\hline
\textbf{Kompetansemål} & \textbf{Læringsmål} & \textbf{Arbeidsmåter} & \textbf{Vurdering} \\
\hline
Naturfag: Å kunne regne i naturfag er å kunne innhente, bearbeide og framstille relevant tallmateriale (...)\newline

Matematikk: lese, hente ut og vurdere matematikk i relevante tekstar om ulike tema og presentere relevante berekningar og analysar av resultata
& 
Elevene skal kunne: loggføre egene innhentede data \newline

sammenligne med historiske data\newline

presentere data i velvalgte diagrammer
& 
Elevene sammenligner pH-målingene fra ekskursjonen med historiske data og naturlige forventede verdier.\newline

De lager egnede diagrammer for å presentere resultatene og diskuterer observasjonene i grupper.
& 
Korrekt analyse\newline

Valg av grafisk framstilling\newline

Evne til å trekke faglige konklusjoner
\\
\hline
\end{tabular}

\newpage

% ------------------ ØKT 3 ------------------
\subsection*{Undervisningsøkt 3 – Debattinnlegg 2-4 timer}

\begin{tabular}{|p{5cm}|p{6cm}|p{8cm}|p{6cm}|}
\hline
\textbf{Kompetansemål} & \textbf{Læringsmål} & \textbf{Arbeidsmåter} & \textbf{Vurdering} \\
\hline
Norsk: lytte til andre, bygge opp saklig argumentasjon og bruke retoriske appellformer i diskusjoner\newline

bruke ulike kilder på en kritisk, selvstendig og etterrettelig måte\newline

kombinere virkemidler og uttrykksformer kreativt i egen tekstskaping\newline

I norsk handler det tverrfaglige temaet folkehelse og livsmestring om å utvikle elevenes evne til å uttrykke seg skriftlig og muntlig\newline

Naturfag: utforske en selvvalgt naturfaglig problemstilling, presentere funn og argumentere for valg av metoder\newline

I naturfag handler det tverrfaglige temaet bærekraftig utvikling om at elevene skal få kompetanse til å gjøre miljøbevisste valg og handlinger, og se disse i sammenheng med lokale og globale miljø- og klimautfordringer (...)
& 
Elevene skal kunne: skrive et debattinnlegg basert på egne data og refleksjoner\newline

bruke kilder på en forsvalig måte, og presentere data fra disse på en ærlig måte\newline

argumentere saklig 
& 
Elevene bruker resultatene fra tidligere økter til å skrive et debattinnlegg. Tema velges fritt, f.eks.: Var kalking av vassdrag et nødvendig inngrep? Er det behov for kalking i dag? etc.
& 
Faglig korrekt argumentasjon\newline

Bruk av data fra ekskursjon og analyse\newline

Klar og strukturert formidling
\\
\hline
\end{tabular}

\end{landscape}

\end{document}